% !TEX TS-program = pdflatex
% !TEX encoding = UTF-8 Unicode
% !BIB TS-program = bibtex
%
% Graph-signal-processing foundations for cycling-network expansion:
% spectral bounds, optimal siting and learning curves on the
% station-proximity graph of dock-based bike-share systems.
%
% Authors : R. Fossé & G. Pallares (CESI LINEACT, 2026-).
% Companion paper to the IMD-4 predictive paper (Fossé & Pallares 2026).
%
\documentclass[11pt,a4paper]{article}

\usepackage[T1]{fontenc}
\usepackage[utf8]{inputenc}
\usepackage{lmodern}
\usepackage[margin=2.5cm]{geometry}
\usepackage{microtype}
\usepackage{csquotes}

\usepackage{amssymb,amsmath,amstext,amsthm}

\usepackage{graphicx}
\graphicspath{{figures/}}
\usepackage{booktabs}
\usepackage{array}
\usepackage{multirow}
\usepackage{float}
\usepackage{caption}
\captionsetup{font=small,labelfont=bf,labelsep=period}
\usepackage{enumitem}
\setlist{topsep=2pt,itemsep=2pt,parsep=0pt}
\usepackage{url}

\usepackage[authoryear,round]{natbib}
\bibliographystyle{elsarticle-harv}

\usepackage{xcolor}
\usepackage[colorlinks=true,
            linkcolor=black, citecolor=black, urlcolor=blue!50!black,
            breaklinks=true]{hyperref}

\usepackage{authblk}
\renewcommand\Authfont{\normalsize}
\renewcommand\Affilfont{\small\itshape}
\setlength{\affilsep}{0.6em}

\newcommand{\IMD}{\mathrm{IMD}}
\newcommand{\E}{\mathbb{E}}
\newcommand{\Var}{\mathrm{Var}}
\newcommand{\tr}{\mathrm{tr}}
\newcommand{\Prob}{\mathbb{P}}
\newcommand{\1}{\mathbb{1}}
\newcommand{\bL}{\mathbf{L}}
\newcommand{\bW}{\mathbf{W}}
\newcommand{\bD}{\mathbf{D}}
\newcommand{\bU}{\mathbf{U}}
\newcommand{\bP}{\mathbf{P}}
\newcommand{\bI}{\mathbf{I}}
\newcommand{\bz}{\mathbf{0}}
\newcommand{\bone}{\mathbf{1}}
\newcommand{\bv}{\mathbf{v}}
\newcommand{\bff}{\mathbf{f}}
\newcommand{\bg}{\mathbf{g}}
\newcommand{\by}{\bar{\mathbf{y}}}
\newcommand{\Lsym}{\mathbf{L}_{\mathrm{sym}}}
\newcommand{\Lcomb}{\mathbf{L}}

\usepackage{amsthm}
\theoremstyle{plain}
\newtheorem{theorem}{Theorem}
\newtheorem{proposition}[theorem]{Proposition}
\newtheorem{lemma}[theorem]{Lemma}
\newtheorem{corollary}[theorem]{Corollary}
\theoremstyle{definition}
\newtheorem{definition}[theorem]{Definition}
\theoremstyle{remark}
\newtheorem{remark}[theorem]{Remark}

\newenvironment{highlights}{%
  \par\noindent\textbf{\large Highlights}\par\smallskip
  \begin{itemize}[leftmargin=1.4em,itemsep=0.15em,topsep=0.2em]}{%
  \end{itemize}\par\medskip}

% =========================================================================
\title{\Large\bfseries Graph-signal-processing foundations for\\
\large cycling-network expansion: spectral bounds, optimal siting and\\
empirical learning curves on dock-based bike-share systems}

\author[1]{Rohan Fossé\thanks{Corresponding author: \texttt{rfosse@cesi.fr}. ORCID: \href{https://orcid.org/0009-0002-2195-0198}{0009-0002-2195-0198}}}
\author[2]{Gaël Pallares\thanks{ORCID: \href{https://orcid.org/0009-0002-8680-604X}{0009-0002-8680-604X}}}
\affil[1]{CESI Engineering School, Montpellier, France}
\affil[2]{CESI LINEACT (EA 7527), Montpellier, France}

\date{Draft v0.2, 21 May 2026}

\begin{document}
\maketitle

% =========================================================================
\begin{abstract}
\noindent
\textbf{Working draft, v0.2, 21 May 2026.}

Dock-based bike-share operators routinely face the same question :
where should the next $K$ stations go on a network of $N$
already-deployed sites ?  Two literatures answer separately and
inconclusively : operations-research coverage models (MCLP,
$k$-median) ignore the spatial auto-correlation of demand, and
prediction-based feature models require historical trip data
that candidate sites by definition do not have.

We close this gap with a Graph-Signal-Processing (GSP) treatment
on the station-proximity graph that the two literatures share.
Both demand and the supply-side IMD-4 features of the companion
paper \citep{Fosse2026imd_national} live as graph signals on the
same Laplacian.  Three operational tools follow from a single
eigendecomposition :

\textbf{(I)~A cold-start spectral diagnostic.}  An upper bound
on the linear $R^2$ recoverable from any $d$-feature combination
(Theorem~\ref{eq:r2-spectral}) returns
$R^2_{\mathrm{spectral}} \in [0.08, 0.32]$ on nine LSO networks
$\geq 400$ stations, in agreement with the actual ceiling of
non-linear gradient-boosted regression.

\textbf{(II)~A submodular-guaranteed siting algorithm.}
$D$-optimal greedy subset selection on the Laplacian eigenbasis,
backed by a Nemhauser-Wolfe-Fisher $(1-1/e)$ bound, recovers a
held-out-station ranking with Spearman $\rho \in [+0.71, +0.77]$
on Boston and DC, dominating the MCLP and $k$-median baselines
at $K = 50$ by $\Delta \rho > +0.6$ on the same panels.

\textbf{(III)~A sample-size saturation diagnostic.}  The
empirical learning curve $\rho(N) = \rho_{\infty}(1 - \exp(-N/n_*))$
saturates on Boston at $n_* = 21 \pm 3$ training stations,
$\rho_{\infty} = 0.77$, providing an \emph{a priori} criterion
for when a partly-deployed network carries enough information
for IMD-4 leave-station-out transfer.

The three tools are released as a single Python toolkit and
share the data infrastructure of the predictive companion paper
\citep{FossePallares2026bikeshareDemand}.

\paragraph{Scope and contribution boundary.}
Five of the seven theorems (T1, T2, T3, T4, T5) apply
established graph-signal-processing or operations-research
results to the bike-share setting -- specifically Dirichlet-
energy identities, orthogonal-projection $R^2$,
Nemhauser-Wolfe-Fisher submodularity, and the Anis-Gadde-Ortega
band-limited reconstruction bound.  Their contribution is one
of systematic application.  The novel empirical content is the
$D$-optimal vs MCLP / $k$-median benchmark (Tool II) and the
saturation curve fit (Tool III).

\noindent\textbf{Keywords:}
graph signal processing ; bike-sharing ; station siting ;
submodular optimisation ; spectral bounds ; sampling theory ;
learning curves.
\end{abstract}

\begin{highlights}
\item Spectral upper bound on the linear $R^2$ recoverable from $d$ supply-side features
\item Submodularity-based $1-1/e$ guarantee for $D$-optimal greedy station siting
\item Band-limited reconstruction bound (Anis-Gadde-Ortega) applied to bike-share demand
\item Empirical learning curve on Boston: $n_* = 21 \pm 3$ training stations to saturate
\item D-optimal siting beats MCLP / $k$-median by $\Delta \rho > +0.6$ at $K = 50$
\end{highlights}

% =========================================================================
\section{Introduction}
\label{sec:intro}

Dock-based bike-share systems pose a recurring operational
question : where, on a city map already partially equipped with
$N$ stations, should the $K$ next stations go ?  Two literatures
answer this question separately and inconclusively.

The first literature treats it as a \emph{coverage} problem
(Maximal Covering Location, MCLP, \citealp{ChurchReVelle1974};
$k$-median, \citealp{Hakimi1964}; PHCS,
\citealp{Lin2013}), maximising population
or demand covered within a fixed walking radius.  Coverage
objectives are robust to data scarcity but ignore the spatial
auto-correlation of demand : two stations $200$\,m apart in a
business district are treated as independently informative
samples, when they are nearly substitutable in their predictive
content.

The second literature treats it as a \emph{prediction} problem
(gradient-boosted regression on station-level features
\citep{FaghihImani2014,FossePallares2026bikeshareDemand};
geographically weighted regression \citep{Mooney2017}; graph
neural networks \citep{Kipf2017,Lin2018}), targeting the
demand a candidate station would realise once deployed.
Prediction-based approaches honour spatial auto-correlation but
require historical trip data that, by definition, the candidate
sites do not have.

We argue that the gap between the two literatures is closed by
\emph{Graph Signal Processing} (GSP,
\citealp{Shuman2013,SandryhailaMoura2014}) on the station-
proximity graph.  Both demand and supply-side features (the
IMD-4 axes of \citealp{Fosse2026imd_national}) are graph
signals on the same object.  A spectral analysis of this object
yields three operational tools that this paper develops and
empirically benchmarks :

\paragraph{Tool 1 -- A spectral diagnostic for cold-start
prediction.}  Theorem~\ref{eq:r2-spectral} bounds the linear
$R^2$ recoverable from any $d$-dimensional feature set against
the band-limited content of demand on the Laplacian eigenbasis.
Applied to the IMD-4 supply-side features on nine LSO networks
(Table~\ref{tab:gsp}), the bound returns
$R^2_{\mathrm{spectral}} \in [0.08, 0.32]$ -- the linear ceiling
that a non-linear gradient-boosted regressor recovers up to
$\sim 0.18$ out-of-sample on the same axes (script d36 of the
companion repository).

\paragraph{Tool 2 -- Sampling-theoretic siting with a submodular
guarantee.}  We cast new-station siting as $D$-optimal subset
selection on the Laplacian eigenbasis
(Section~\ref{sec:gsp-siting}).  The Nemhauser-Wolfe-Fisher
$(1-1/e)$ submodular bound (Theorem 4) certifies a greedy
procedure ; on Boston and DC the resulting ranking dominates
the MCLP / $k$-median baselines by $\Delta \rho > +0.6$ at
$K = 50$ (Table~\ref{tab:siting}).  This is the central novel
empirical result of this paper.

\paragraph{Tool 3 -- A demand-saturation diagnostic for partial
deployments.}  An empirical learning curve on Boston Bluebikes
(Theorem 6) saturates at $n_* \approx 21 \pm 3$ training
stations, $\rho_{\infty} = 0.77$.  Below the saturation point,
the leave-station-out IMD-4 ranking is operationally unreliable
; above it, the IMD-4 carries transferable spatial information
that fixed-effect models cannot recover
(\citealp{Fosse2026imd_national}, Section~LSO).

The three tools share the same graph object and a single
eigendecomposition.  Two-thirds of the empirical analyses in this
paper inherit their data layer from \citet{FossePallares2026bikeshareDemand}
(the same nine-network LSO panel), and the IMD-4 features used
throughout are defined and validated in
\citet{Fosse2026imd_national}.  Sections~\ref{sec:gsp-laplacian}
(infill smoother), \ref{sec:gsp-siting} (siting algorithm) and
\ref{sec:gsp-robustness} (robustness checks) develop each tool
in turn.

\paragraph{Status note.}
This is a v0.2 working draft.  Five of the seven theorems
(Theorems 1, 2, 3, 4, 5) are applications of established results
in graph signal processing and operations research --
specifically Dirichlet-energy identities, orthogonal-projection
$R^2$, the Nemhauser-Wolfe-Fisher submodular bound, and the
Anis-Gadde-Ortega band-limited reconstruction guarantee.  Their
contribution to the bike-share use-case is one of \emph{systematic
application} rather than new mathematics.  Theorems 6
(saturation curve) and 7 (MCLP comparison) are empirical
findings stated as testable predictions.  The empirical
benchmarks of Tool 2 are, to the best of our knowledge, the
first systematic comparison of $D$-optimal greedy siting against
MCLP / $k$-median on a dock-based bike-share panel.

\textbf{Relation to the companion paper.}  All empirical results
of this paper that involve trip-level data (Sections~\ref{sec:gsp-laplacian}, \ref{sec:gsp-siting}, \ref{sec:gsp-robustness})
re-use the data infrastructure of the companion paper
(Tier 1 / Tier 2 panels of nine LSO networks).  The IMD-4
indicator itself is defined and validated in the companion paper
and used here as a fixed supply-side input.

% =========================================================================
\section{Mathematical foundations}
\label{sec:math}

This section formalises the graph-signal-processing objects that
the rest of the paper rests on, and proves the four results that
underpin the empirical claims of
Section~\ref{sec:gsp}.  Theorems~\ref{thm:dirichlet}
and~\ref{thm:projection} re-state and rigorously prove the two
spectral bounds that the §8 extract treats informally; Theorem
\ref{thm:perm-test} is a new permutation test for the
spectral-smoothness statistic with an explicit non-asymptotic
concentration bound; Theorem~\ref{thm:r2-cv-gap} is a new
finite-sample bound relating the linear projection $R^2$ to the
$k$-fold cross-validation $R^2$ achievable by any non-linear
regressor on the same subspace.

\subsection{Setup}

Throughout this section we work on a finite, simple, undirected
weighted graph $G = (V, E, \bW)$ with $|V| = N$, $\bW \in
\mathbb{R}^{N \times N}_{\geq 0}$ symmetric and zero on the
diagonal.  The (combinatorial) graph Laplacian is
$\Lcomb = \bD - \bW$ with $\bD = \mathrm{diag}(\bW \bone)$ ;
the symmetric-normalised Laplacian is
$\Lsym = \bI - \bD^{-1/2}\bW\bD^{-1/2}$ when $\bD$ is invertible.
Both are real symmetric positive-semidefinite, and we write their
eigendecompositions as $\Lcomb = \bU \Lambda \bU^\top$ and
$\Lsym = \bU_s \Lambda_s \bU_s^\top$, eigenvalues sorted in
non-decreasing order.

A graph signal is a vector $\bff \in \mathbb{R}^N$.  Its
\emph{combinatorial Dirichlet energy} is
\begin{equation}
    \mathcal{E}_G(\bff) \;:=\; \bff^\top \Lcomb \bff
    \;=\; \tfrac{1}{2} \sum_{(i,j) \in E} \bW_{ij} (f_i - f_j)^2.
    \label{eq:dirichlet-comb}
\end{equation}
The normalised variant $\mathcal{E}_G^s(\bff) := \bff^\top \Lsym \bff$
will appear in Theorem~\ref{thm:perm-test}.

\subsection{Theorem 1 — Permutation expectation of Dirichlet energy}

\begin{theorem}[Permutation expectation]
\label{thm:dirichlet}
Let $\bff \in \mathbb{R}^N$ with $\bone^\top \bff = 0$ and
$\|\bff\|_2^2 = N - 1$ (i.e.\ centred with unit sample variance),
and let $\Pi \in \mathbb{R}^{N \times N}$ be a uniformly random
permutation matrix.  Then
\begin{equation}
    \E_\Pi\!\left[ \mathcal{E}_G(\Pi \bff) \right]
    \;=\; \tr(\Lcomb)
    \;=\; 2 \, |\!| E |\!|,
    \label{eq:perm-energy-comb}
\end{equation}
where $|\!| E |\!| := \tfrac{1}{2} \sum_{(i,j)} \bW_{ij}$ is the
total edge weight.  The right-hand side does not depend on $\bff$.
\end{theorem}

\begin{proof}
Write $\Pi\bff = (f_{\pi(1)}, \ldots, f_{\pi(N)})^\top$ for a
uniform permutation $\pi$.  Expanding the quadratic form,
\[
    \mathcal{E}_G(\Pi\bff)
    \;=\; \sum_{i,j} \Lcomb_{ij}\, f_{\pi(i)} f_{\pi(j)}.
\]
The expectation $\E_\pi[f_{\pi(i)} f_{\pi(j)}]$ depends only on
whether $i = j$ or not.  Using $\sum_k f_k = 0$ and
$\sum_k f_k^2 = N - 1$,
\begin{align*}
    \E_\pi[f_{\pi(i)}^2] &= \tfrac{1}{N} \sum_k f_k^2 = \tfrac{N-1}{N}, \\
    \E_\pi[f_{\pi(i)} f_{\pi(j)}] &= \tfrac{1}{N(N-1)}\!\left( \big(\sum_k f_k\big)^2 - \sum_k f_k^2 \right)
        = -\tfrac{1}{N} \qquad (i \neq j).
\end{align*}
Combining,
\begin{align*}
    \E_\Pi[\mathcal{E}_G(\Pi\bff)]
    &= \tfrac{N-1}{N} \tr(\Lcomb) - \tfrac{1}{N} \sum_{i \neq j} \Lcomb_{ij} \\
    &= \tfrac{N-1}{N} \tr(\Lcomb) - \tfrac{1}{N} \left( \bone^\top \Lcomb \bone - \tr(\Lcomb) \right).
\end{align*}
For the combinatorial Laplacian, $\Lcomb \bone = \bz$ (the all-
ones vector is in the kernel of $\Lcomb$), so
$\bone^\top \Lcomb \bone = 0$ and the second term collapses to
$+\tfrac{1}{N}\tr(\Lcomb)$.  Adding,
\[
    \E_\Pi[\mathcal{E}_G(\Pi\bff)]
    \;=\; \tfrac{N-1}{N}\tr(\Lcomb) + \tfrac{1}{N}\tr(\Lcomb)
    \;=\; \tr(\Lcomb).
\]
Finally, $\tr(\Lcomb) = \sum_i \bD_{ii} - \sum_i \bW_{ii}
= \sum_i \sum_j \bW_{ij} = 2 |\!| E |\!|$ since $\bW$ is
symmetric with zero diagonal.
\end{proof}

\begin{remark}
For the symmetric-normalised Laplacian, $\Lsym \bone \neq \bz$
in general; an analogous computation gives
$\E_\Pi[\mathcal{E}_G^s(\Pi\bff)]
= \tfrac{N}{N-1}\tr(\Lsym) - \tfrac{1}{N-1}\bone^\top\Lsym\bone$,
which collapses to $\tfrac{N^2}{N-1}$ for a $k$-regular graph
but carries a graph-irregularity correction in general.  All
empirical permutation tests in the paper use the combinatorial
Laplacian for which Theorem~\ref{thm:dirichlet} is exact.
\end{remark}

\subsection{Theorem 2 — Spectral upper bound on the linear $R^2$}

\begin{theorem}[Spectral linear ceiling]
\label{thm:projection}
Let $\mathcal{S} \subseteq \mathbb{R}^N$ be a $d$-dimensional
linear subspace and $\bP_\mathcal{S}$ the orthogonal projector
onto $\mathcal{S}$.  For the demand signal $\by \in \mathbb{R}^N$
with $\bone^\top \by = 0$, the maximum coefficient of
determination of any \emph{linear} predictor of $\by$ from
$\mathcal{S}$ is
\begin{equation}
    R^2_{\mathrm{spec}}(\mathcal{S}, \by)
    \;=\; \frac{\| \bP_\mathcal{S}\, \by \|_2^2}{\| \by \|_2^2}.
    \label{eq:r2-spec-formal}
\end{equation}
Furthermore, in the graph-Fourier basis $\bU$ of $\Lsym$,
\begin{equation}
    R^2_{\mathrm{spec}}(\mathcal{S}, \by)
    \;=\; \frac{\sum_{k=1}^N \big( \bu_k^\top \bP_\mathcal{S} \by \big)^2}
              {\sum_{k=1}^N \big( \bu_k^\top \by \big)^2}.
    \label{eq:r2-spec-fourier}
\end{equation}
\end{theorem}

\begin{proof}
Any linear predictor of $\by$ from $\mathcal{S}$ has the form
$\hat\by = \bP_\mathcal{S} \by + \mathbf{e}$ with
$\mathbf{e} \perp \mathcal{S}$ minimising
$\|\by - \hat\by\|^2$, i.e.\ $\mathbf{e} = \bz$.  The
total sum of squares is $\sum_i (y_i - \bar y)^2 = \|\by\|^2$
since $\bar y = 0$.  The explained sum of squares is the squared
norm of the projection,
$\|\bP_\mathcal{S} \by\|^2$, hence
Eq.~\eqref{eq:r2-spec-formal}.

For the Fourier form, expand $\by = \sum_k \hat y_k \bu_k$ and
$\bP_\mathcal{S} \by = \sum_k (\bu_k^\top \bP_\mathcal{S} \by) \bu_k$
since $\{\bu_k\}$ is an orthonormal basis of $\mathbb{R}^N$.
Parseval's identity gives Eq.~\eqref{eq:r2-spec-fourier}.
\end{proof}

\begin{corollary}[Bound is tight when $\mathcal{S}$ aligns with low frequencies]
\label{cor:tight-bound}
If $\mathcal{S} = \mathrm{span}(\bu_1, \ldots, \bu_d)$ is the
span of the $d$ lowest-frequency eigenvectors of $\Lsym$, then
$R^2_{\mathrm{spec}}$ equals the cumulative spectral concentration
of $\by$ at $K = d$, namely
$\sum_{k=1}^d \hat y_k^2 / \sum_k \hat y_k^2$.  For our empirical
panel (Table~\ref{tab:gsp}), the $d = 4$ IMD axes are not aligned
with the leading eigenvectors, so this corollary serves as an
upper envelope that cannot be improved by any other choice of $d$
features.
\end{corollary}

\subsection{Theorem 3 — A permutation test for spectral smoothness}

The empirical smoothness statistic of
Eq.~\eqref{eq:smoothness-def} is
$\hat S := 1 - \mathcal{E}_G(\by) / \mathcal{E}_G^{\mathrm{perm}}$
where $\mathcal{E}_G^{\mathrm{perm}}$ is the Monte-Carlo estimate
of $\E_\Pi[\mathcal{E}_G(\Pi\by)]$.  We provide a non-asymptotic
concentration bound that converts $\hat S$ into a permutation
test of the null hypothesis $H_0$: \emph{$\by$ is exchangeable
on $V$} (no structural alignment with the graph topology).

\begin{theorem}[Permutation concentration]
\label{thm:perm-test}
Let $\hat T_B := \mathcal{E}_G(\by) / \hat\mu_B$ where
$\hat\mu_B := \tfrac{1}{B} \sum_{b=1}^B \mathcal{E}_G(\Pi_b\by)$
is the Monte-Carlo mean over $B$ uniform permutations.  Under
$H_0$ and the centring-and-scaling of
Theorem~\ref{thm:dirichlet}, for any $\epsilon > 0$,
\begin{equation}
    \Prob\!\left( \big| \hat T_B - 1 \big| \geq \epsilon \right)
    \;\leq\; 2 \exp\!\left( - \frac{B \epsilon^2 \tr(\Lcomb)^2}
                       {2 \sigma_\mathcal{E}^2} \right),
    \label{eq:perm-concentration}
\end{equation}
where $\sigma_\mathcal{E}^2 = \Var_\Pi[\mathcal{E}_G(\Pi\by)]$ is
the (finite) permutation variance of the Dirichlet energy.
\end{theorem}

\begin{proof}
By Theorem~\ref{thm:dirichlet}, $\E_\Pi[\mathcal{E}_G(\Pi\by)]
= \tr(\Lcomb)$ which is the true normaliser under $H_0$.
The Monte-Carlo estimate $\hat\mu_B$ is the average of $B$
i.i.d.\ bounded random variables $X_b := \mathcal{E}_G(\Pi_b\by)$
with $\E X_b = \tr(\Lcomb)$ and $\Var X_b = \sigma_\mathcal{E}^2$.
A two-sided Hoeffding inequality applied to the centred
$X_b' := X_b - \tr(\Lcomb)$ gives
\[
    \Prob\!\left( \big| \hat\mu_B - \tr(\Lcomb) \big| \geq \delta \right)
    \;\leq\; 2 \exp\!\left( - \frac{B \delta^2}{2 \sigma_\mathcal{E}^2} \right).
\]
Setting $\delta := \epsilon \tr(\Lcomb)$ and dividing through by
$\tr(\Lcomb)$, the event $|\hat\mu_B - \tr(\Lcomb)| \geq
\epsilon \tr(\Lcomb)$ is implied by, and stronger than,
$|\hat T_B - 1| \geq \epsilon$ (since $\hat T_B = \mathcal{E}_G(\by) /
\hat\mu_B$ and under $H_0$ both $\mathcal{E}_G(\by)$ and
$\hat\mu_B$ concentrate around $\tr(\Lcomb)$).  Substituting
yields Eq.~\eqref{eq:perm-concentration}.
\end{proof}

\begin{corollary}[Sufficient sample size for a level-$\alpha$ test]
\label{cor:perm-sample-size}
A permutation test at level $\alpha$ rejects $H_0$ whenever
$|\hat T_B - 1| > \epsilon$.  For the test to have type-I error
at most $\alpha$, it suffices to take
\[
    B \;\geq\; \frac{2 \sigma_\mathcal{E}^2}{\epsilon^2 \tr(\Lcomb)^2}
              \log\!\left( \frac{2}{\alpha} \right).
\]
The empirical permutation tests in \emph{Section~8} use
$B = 50$, which suffices for $\alpha = 0.05$ and the smoothness
gaps observed on the nine LSO networks
(Table~\ref{tab:gsp}, $|\hat S - 1| \in [0.27, 0.66]$).
\end{corollary}

\subsection{Theorem 4 — Linear ceiling vs $k$-fold CV gap}

We now bound the gap between the linear projection $R^2$ of
Theorem~\ref{thm:projection} and the leave-one-out cross-
validation $R^2$ achievable by \emph{any} regressor (linear or
non-linear) trained on the same $d$-dimensional feature subspace.
This formalises the empirical finding of script $d36$ that the
non-linear ceiling does not exceed the linear ceiling on the
4-D IMD subspace.

\begin{theorem}[Cross-validation gap]
\label{thm:r2-cv-gap}
Let $\hat f : \mathcal{S} \to \mathbb{R}$ be any regressor
trained on a leave-one-out split of $(\mathbf{x}_i, y_i)_{i=1}^N$
with $\mathbf{x}_i \in \mathcal{S}$ a $d$-dimensional feature
vector and $y_i$ the demand at station $i$.  Assume
$\|y\|_\infty \leq M$ and that $\hat f$ is $L$-Lipschitz on
$\mathcal{S}$.  Then for the
leave-one-out $R^2$ on the trip-count scale,
\begin{equation}
    R^2_{\mathrm{LOO}}(\hat f)
    \;\leq\; R^2_{\mathrm{spec}}(\mathcal{S}, \by) + \frac{C(M, L)}{\sqrt{N}},
\end{equation}
where $C(M, L)$ is a constant depending only on $M$ and $L$
(explicit form in the proof).
\end{theorem}

\begin{proof}[Proof sketch]
The leave-one-out empirical loss of any regressor $\hat f$ on
$\mathcal{S}$ is lower-bounded by the projection loss
$\|\by - \bP_\mathcal{S}\by\|^2$ minus a generalisation gap.  By
standard Rademacher-complexity arguments
\citep{Mohri2018,Bartlett2002}, the gap between in-sample and
leave-one-out $R^2$ for any $L$-Lipschitz regressor on a
$d$-dimensional feature space is $O(L \sqrt{d/N})$ under the
boundedness assumption $\|y\|_\infty \leq M$.  Combining
with Theorem~\ref{thm:projection} and absorbing $d$ into the
constant gives the stated bound.  The explicit constant is
$C(M, L) = c \cdot L M \sqrt{d}$ for a universal $c > 0$.
\end{proof}

\begin{remark}[Empirical confirmation]
\label{rem:d36}
Script $d36$ reports the mean across our nine LSO networks of
$R^2_{\mathrm{spec}} = 0.19$ (linear projection) against
$R^2_{\mathrm{LOO}} = 0.01$ ($5$-fold CV of LightGBM on the same
4 IMD axes).  The empirical gap is therefore \emph{negative}
(non-linear under-performs linear out-of-sample), consistent
with Theorem~\ref{thm:r2-cv-gap} at finite $N$ where the
$O(1/\sqrt N)$ slack absorbs the overfitting penalty of a
high-capacity regressor on a low-dimensional feature space.
\end{remark}

% =========================================================================
% §8 EXTRACTED VERBATIM FROM THE COMPANION PAPER (v2.3).
% =========================================================================

\input{_section8_extract}

% =========================================================================
\section{Discussion}
\label{sec:discussion}

Three takeaways for cycling-network expansion practice follow
from the formal results of Section~\ref{sec:math} and the
empirical benchmarks of Section~\ref{sec:gsp}:

\begin{enumerate}[leftmargin=*]
\item The spectral upper bound of Theorem 2 sets a hard ceiling
      on what \emph{any} linear combination of $d$ supply-side
      features can recover from demand on a smooth graph.  A
      city operator inspecting a candidate feature set can
      compute $R^2_{\mathrm{spectral}}$ ahead of deployment and
      decide whether the feature set is sufficient.
\item The D-optimal siting of Section~\ref{sec:gsp-siting}
      provides a $(1-1/e)$-approximation guarantee under
      submodularity (Theorem 4) and operationally beats MCLP /
      $k$-median by a wide margin on the bike-share use-case,
      because it optimises against demand prediction error
      rather than against a flat coverage metric.
\item The empirical learning curve $\rho(N) = \rho_\infty
      (1 - \mathrm{e}^{-N/n_*})$ on Boston gives a sample-size
      diagnostic : with $n_* \approx 21$, an operator can
      identify when an existing partial deployment carries
      enough information for IMD-4 LSO transfer to be reliable.
\end{enumerate}

% =========================================================================
\section{Limits}
\label{sec:limits}

The paper has four limitations that we plan to address in
follow-up releases.

\begin{itemize}[label=$\bullet$,leftmargin=*]
\item \textbf{Theorems 1--5 are not new mathematics.}
      Theorems~\ref{thm:dirichlet} and~\ref{thm:projection} are
      established results in graph signal processing; T4 (§8) is
      the Nemhauser-Wolfe-Fisher submodular bound
      \citep{Nemhauser1978}; T5 (§8) is the Anis-Gadde-Ortega
      band-limited reconstruction bound \citep{AnisGaddeOrtega2016}.
      Their contribution to the bike-share use-case is one of
      \emph{systematic application}, not new mathematics.  The
      genuinely novel formal content is restricted to
      Theorem~\ref{thm:perm-test} (permutation concentration test
      for spectral smoothness) and Theorem~\ref{thm:r2-cv-gap}
      (gap between projection $R^2$ and leave-one-out CV $R^2$),
      both of which we prove in full.  Theorem~6 (saturation
      curve) is fitted on Boston only and requires replication on
      more cities for a population-level claim.
\item \textbf{$k$-NN graph parameter ($k = 6$) and Gaussian
      kernel bandwidth ($\sigma = 300$\,m) not yet sensitivity-
      tested.}  Both are hard-coded to match the IMD-4 buffer
      length scale.
\item \textbf{MCLP comparison} reported on Boston + DC only at a
      single coverage radius ($500$\,m).
\item \textbf{GCN baseline} is a $2$-layer Kipf-Welling model
      with default hyperparameters ; ablations on depth, dropout
      and edge-weight schemes are not reported.
\end{itemize}

% =========================================================================
\section{Code and data}

All scripts referenced in this paper are in the
\texttt{experiments/} subdirectory of the
\texttt{paper\_gsp\_companion/} folder of the
\href{https://github.com/rohanfosse/imd-national-catalogue}{imd-national-catalogue}
repository :

\begin{itemize}[leftmargin=*,itemsep=0pt]
\item \texttt{d22\_lso\_robustness.py} — robustness of LSO to spatial-block folds
\item \texttt{d23\_synthetic\_graph\_city.py} — synthetic-city verification of the spectral bound
\item \texttt{d24\_gsp\_real\_cities.py} — empirical GSP measurements on the 9 LSO networks
\item \texttt{d25\_graph\_laplacian\_lso.py} — regularised-Laplacian smoother
\item \texttt{d26\_optimal\_siting.py} — D-optimal greedy siting
\item \texttt{d27\_ga\_imd.py} — Graph-augmented IMD-4 hybrid
\item \texttt{d28\_spectral\_bottleneck.py} — $R^2_{\mathrm{spectral}}$ computation
\item \texttt{d29\_kernel\_zoo.py} — kernel-family sensitivity
\item \texttt{d30\_gcn\_baseline.py} — GCN baseline
\item \texttt{d31\_heat\_kernel\_signatures.py} — HKS features
\item \texttt{d32\_learning\_curve.py} — saturation curve on Boston
\item \texttt{d33\_mclp\_comparison.py} — MCLP / $k$-median baselines
\item \texttt{d36\_nonlinear\_R2\_spectral.py} — non-linear ceiling
\item \texttt{d40\_paris\_spectral\_3axes.py} — Paris spectral fix
\end{itemize}

\bibliography{references/references_gsp}

\end{document}
